\documentclass{handout}

\assignment{Lab 03}
\title{Motion in 2D}

\begin{document}
\maketitle

\section*{Introduction}
We have studied free fall in one and two dimensions in class. Let's get more of a feel for motion graphs and why they help represent a schematic `movie' of motion. We will do this by analyzing a video one frame at a time.


\section*{Planning a suitable video}
The video analysis software we will be using assumes two things:
\begin{enumerate}
    \item All motion occurs perpendicular to the camera's line of sight
    \item The object being tracked has a single position that can be indicated every frame.
    \item An object of known size is included in the frame.
\end{enumerate}

\begin{questions}
\question You need a location to film your movie: where would you like to go where you can reliably throw an object > 1 story with a clear line of sight?
\answerspace{2in}

\question Choose a ball that will be relatively easy to track. Which ball are you choosing and why? 

\answerspace{2in}

\question If the camera is pointed east, what directions of motion can the video detect, and why?
\begin{parts}
\part North
\part South
\part East 
\part West
\part Up
\part Down
\end{parts}

\question How do you plan to mount or hold your camera? The camera frame should be fixed.

\answerspace{2in}

\end{questions}

\section*{Record your video}
Now let’s go record your videos. Make two videos, both \textbf{including the object landing}:
\begin{enumerate}
    \item Throwing your object straight up.
    \item Tossing your object forward.
\end{enumerate}
Weather permitting, we will go outside, since daylight is generally brighter
than indoor lighting. The shorter the video the better, so don’t start recording
until you’re ready to throw, and go for just a single throw per video. Try not
to throw the object out of the camera field of view. This will take some trial
and error. Good luck, and let me know if you need any help.

Once you are done, somehow get the videos onto the lab computer. Make sure to  log into the lab computer in incognito mode (CTRL+SHIFT+N in the browser). This way you don't stay logged into the computer.

\section*{Video tracking}
Now, let’s analyze your videos. We will do this with software called ``PASCO Capstone''.
\textbf{You must analyze both of your videos.}

Tracker is installed on the desktop computers in the lab. Find a computer you
like, and copy your video onto the desktop. Your guess is as good as mine for
how best to do this.

Open up the projectile motion lab (brightspace) and follow these steps. 
\begin{enumerate}
    \item On the video pane, open your first video file. 
    \item Drag the reference size to your reference object. This sets the length of the reference object and the origin of your coordinate system.
    \item Make sure the `video analysis mode' is selected in the top left corner of the video pane. You should see `object 1' in the toolbar somewhere.
    \item Use the >> button at the bottom until the first frame of projectile motion.
    \item Click on the location of your object. A red x should appear and the video should advance by one frame.
    \item Continue this until the object hits the ground. Include at least two frames after the object hits the ground.
    \item Save your work just in case.
\end{enumerate}

\begin{questions}
    \question Explain how the motion graphs you see are related to the frames of the video.
    \answerspace{2in}

    \question Explain why the velocity graph looks worse than position, and acceleration looks worse than the velocity graph. 

    \answerspace{2in}

    \question Explain the behavior of the $v_y$ vs $v_x$ graph in terms of what was discussed in class.
    \answerspace{2in}
\end{questions}

For each of the following questions, include a sketch of the motion graph that you used to answer the question. Also indicate where the quantity asked is located on your sketch.

Answer the following for both videos: up and forward.
\begin{questions}
    \question What is the maximum height reached by your object?
    
    Up: \hspace{\fill} Forward: \hspace{\fill}
    \answerspace{1.5in}

    \question At what time did the object achieve that maximum height?
    
    Up: \hspace{\fill} Forward: \hspace{\fill}
    \answerspace{1.5in}

    \question What is the initial vertical velocity of the object?
    
    Up: \hspace{\fill} Forward: \hspace{\fill}
    \answerspace{1.5in}

    \question How far was the object thrown?
    
    Up: \hspace{\fill} Forward: \hspace{\fill}
    \answerspace{1.5in}

    \question How fast was the object thrown (a speed, in m/s)?
    
    Up: \hspace{\fill} Forward: \hspace{\fill}
    \answerspace{1.5in}

    \question What was the direction was the object thrown (an angle above horizontal)?
    
    Up: \hspace{\fill} Forward: \hspace{\fill}
    \answerspace{1.5in}

\end{questions}

Now consider the following calculations:
\begin{questions}
\question Calculate time to achieve maximum from your initial vertical velocity and $g$.

    Up: \hspace{\fill} Forward: \hspace{\fill}
    \answerspace{1.5in}

\question Calculate the maximum height from the initial velocity and time.

    Up: \hspace{\fill} Forward: \hspace{\fill}
    \answerspace{1.5in}

\question Calculate how far your object was thrown from the initial velocity and time.

    Up: \hspace{\fill} Forward: \hspace{\fill}
    \answerspace{1.5in}

\question Explain why these numbers might disagree with your measured values.
\answerspace{1.5in}
\end{questions}
\end{document}